\documentclass[11pt,a4paper]{article}
\usepackage[utf8]{inputenc}
\usepackage[T1]{fontenc}
\usepackage{amsmath,amsfonts,amssymb}
\usepackage{graphicx}
\usepackage{hyperref}
\usepackage{listings}
\usepackage{color}
\usepackage{booktabs}
\usepackage{float}
\usepackage{geometry}
\geometry{margin=1in}
\definecolor{zenblue}{RGB}{41,121,255}
\definecolor{zengreen}{RGB}{52,199,89}
\hypersetup{colorlinks=true,linkcolor=zenblue,urlcolor=zenblue,citecolor=zenblue}

\title{\textbf{Zen-Pro: A Reasoning-Tuned Language Model Built on Qwen3-8B}\\
\large Technical Report v2025.02}
\author{Antje Worring, Zach Kelling \\ Zen LM Research Team\\
\texttt{research@zenlm.org}}
\date{February 2025}

\begin{document}
\maketitle

\begin{abstract}
We present \textbf{Zen-Pro}, a reasoning-oriented instruction model in the Zen family. Zen-Pro is
\emph{not} a 72-billion-parameter model trained from scratch: it is a fine-tuned derivative of
\textbf{Qwen3-8B} \cite{qwen3report}, the 8.2B-parameter open-weight dense model released by
Alibaba's Qwen team under the Apache-2.0 license. Earlier versions of this report claimed a bespoke
``Zen MoDE'' architecture, a 72B parameter count, and a 4.5-trillion-token pretraining run; those
claims were false and have been removed. Zen-Pro applies reasoning-focused post-training---long
chain-of-thought (CoT) supervised fine-tuning and, where a verifiable reward signal is available,
reinforcement learning on math and code---on top of the released Qwen3-8B weights. This report
documents the inherited 8B architecture honestly, describes the post-training we actually perform,
and defers to the upstream Qwen3 technical report for pretraining and benchmark numbers.
\end{abstract}

\tableofcontents
\newpage

%% ─────────────────────────────────────────────────────────────────────────────
\section{Introduction}

Strong reasoning behavior in open models has been driven less by raw parameter count than by
reasoning-focused post-training: long chain-of-thought supervision and reinforcement learning
against verifiable rewards \cite{lightman2023lets, shao2024deepseekmath}. Zen-Pro applies this style
of post-training to a strong, permissively licensed 8B base model.

\paragraph{Provenance and corrections.} Zen-Pro is a derivative of \textbf{Qwen3-8B}, an
8.2B-parameter dense decoder-only transformer released by the Qwen team at Alibaba Cloud under the
Apache-2.0 license \cite{qwen3report, qwen3hf}. Configuration fingerprinting of the released Zen-Pro
weights matches Qwen3-8B (36 layers, hidden size 4096, 32/8 GQA heads, 151{,}936-token vocabulary).
The previously claimed ``72B dense, 4.5T-token, Zen MoDE'' description does not correspond to the
shipped model and has been removed. Zen-Pro is an 8B model, was not pretrained from scratch by the
Zen LM team, and does not introduce a new architecture.

\paragraph{What Zen-Pro adds.} Our contributions are post-training only:

\begin{itemize}
  \item Long chain-of-thought (CoT) supervised fine-tuning on reasoning trajectories, starting from
        Qwen3-8B.
  \item Optional reinforcement learning from verifiable rewards (RLVR) on math and code tasks where
        correctness can be checked programmatically.
  \item Packaged deployment artifacts (BF16 SafeTensors, GGUF, MLX) identical in interface to any
        Qwen3-8B deployment.
\end{itemize}

Zen-Pro is the reasoning-tuned variant in the Zen family; it shares the same 8B Qwen3 base as the
other Zen language models, differing only in post-training emphasis.

%% ─────────────────────────────────────────────────────────────────────────────
\section{Architecture (Inherited from Qwen3-8B)}

\subsection{Model Hyperparameters}

Zen-Pro inherits the Qwen3-8B architecture unchanged. Table~\ref{tab:arch} reproduces the upstream
hyperparameters \cite{qwen3hf}.

\begin{table}[H]
\centering
\caption{Architecture hyperparameters, inherited unchanged from Qwen3-8B \cite{qwen3hf}.}
\label{tab:arch}
\begin{tabular}{lc}
\toprule
\textbf{Hyperparameter} & \textbf{Value} \\
\midrule
Parameters (total)           & 8.2B \\
Non-embedding parameters     & 6.95B \\
Layers                       & 36 \\
Attention heads (query)      & 32 \\
KV heads (GQA)               & 8 \\
Head dimension               & 128 \\
Hidden dimension             & 4096 \\
FFN intermediate dimension   & 12{,}288 \\
Vocabulary size              & 151{,}936 \\
Context length (native)      & 40{,}960 \\
Context length (with YaRN)   & 131{,}072 \\
Position encoding            & RoPE ($\theta = 1{,}000{,}000$) \\
Activation function          & SiLU (SwiGLU) \\
Normalization                & RMSNorm \\
Tied embeddings              & No \\
\bottomrule
\end{tabular}
\end{table}

\subsection{Grouped Query Attention}

Qwen3-8B uses GQA \cite{ainslie2023gqa} with 32 query heads and 8 KV heads (head dimension 128),
reducing the KV cache footprint by 4$\times$ relative to multi-head attention. The KV cache size is

\begin{equation}
  \text{KV cache size} = 2 \cdot L \cdot n_{\text{kv}} \cdot d_k \cdot S \cdot \text{dtype\_bytes}
\end{equation}

with $L=36$ layers, $n_{\text{kv}}=8$ KV heads, $d_k=128$, and sequence length $S$.

\subsection{Context Length}

Qwen3-8B supports a native context of 40{,}960 tokens, extensible to 131{,}072 tokens via YaRN
scaling \cite{peng2023yarn} as documented upstream \cite{qwen3report}. Zen-Pro inherits this
behavior; we do not perform additional context-extension pretraining and rely on the upstream YaRN
configuration for long-context inference.

%% ─────────────────────────────────────────────────────────────────────────────
\section{Post-Training Methodology}

This section describes the reasoning-focused post-training that the Zen LM team applies on top of the
released Qwen3-8B weights. We do not perform pretraining or context extension at scale; for the
pretraining corpus and procedure, see the upstream Qwen3 technical report \cite{qwen3report}.

\subsection{Long Chain-of-Thought Fine-Tuning}

Starting from Qwen3-8B, we apply a CoT-SFT phase on long-form reasoning trajectories, training the
model to produce explicit step-by-step reasoning before a final answer:

\begin{align}
  \mathcal{L}_{\text{CoT-SFT}} &= -\sum_{t=1}^{T_r} \log p(r_t \mid x, r_{<t})
                                   - \lambda \sum_{t=1}^{T_a} \log p(a_t \mid x, r, a_{<t})
\end{align}

where $r$ is the reasoning chain, $a$ is the final answer, and $\lambda$ down-weights the answer
loss to encourage richer intermediate reasoning. Where trajectories are distilled from a stronger
teacher, we attribute the teacher and do not claim the traces as original pretraining data.

\subsection{Reinforcement Learning from Verifiable Rewards (Optional)}

Where a programmatic correctness signal is available, we apply RLVR \cite{lightman2023lets,
cobbe2021gsm8k} on two domains:

\begin{itemize}
  \item \textbf{Mathematics}: symbolic/numeric equivalence checking of final answers.
  \item \textbf{Code}: unit-test execution with pass@k evaluation.
\end{itemize}

A simple verifiable reward with a length penalty,
\begin{equation}
  r(y, y^*) = \mathbb{1}[\text{verify}(y) = y^*] - \alpha \cdot \frac{|y|}{L_{\max}},
\end{equation}
is optimized with Group Relative Policy Optimization (GRPO) \cite{shao2024deepseekmath}, sampling
$G$ responses per prompt and computing within-group advantages to avoid a separate value network.

%% ─────────────────────────────────────────────────────────────────────────────
\section{Evaluation}

\paragraph{On benchmark numbers.} This report does not contain fabricated head-to-head benchmark
tables. For rigorous, independently reproducible results on the Qwen3 base and instruct models
(MMLU, MMLU-Pro, MATH, GSM8K, GPQA, HumanEval, long-context retrieval, and reasoning suites such as
AIME), we refer to the official Qwen3 technical report \cite{qwen3report} and the Qwen3-8B model
card \cite{qwen3hf}. Note that Qwen3 already ships strong reasoning-tuned instruct checkpoints; any
gains we report for Zen-Pro must be measured against the appropriate Qwen3-8B baseline under a
stated harness, and we report only numbers we have actually measured on the released artifact.

\paragraph{Honest framing of expected gains.} Reasoning-focused SFT and RLVR can improve
chain-of-thought accuracy on math and code relative to a non-reasoning baseline, but the magnitude
depends heavily on data and compute and is not assumed here. We do not claim frontier-tier scores
(such as the previously stated MMLU 89.2 / MATH 84.1 / HumanEval 91.3 / GPQA 71.8) for an 8B
model; those numbers were fabricated and are removed. Realistic expectations for an 8B reasoning
finetune should be calibrated against published Qwen3-8B results.

%% ─────────────────────────────────────────────────────────────────────────────
\section{Inference and Deployment}

Because Zen-Pro is architecturally identical to Qwen3-8B, it runs on any Qwen3-compatible stack
(Hugging Face Transformers, vLLM, llama.cpp via GGUF, MLX). At 8B parameters it fits on a single
24\,GB consumer GPU in BF16 and on commodity hardware in 4-bit quantization---substantially more
accessible than the 70B-class deployment described in earlier (incorrect) versions of this report.
We publish BF16 SafeTensors, GGUF (Q4\_K\_M, Q8\_0), and MLX builds. We do not report fabricated
throughput tables; measured throughput matches an equivalent Qwen3-8B deployment under the same
runtime and quantization.

%% ─────────────────────────────────────────────────────────────────────────────
\section{Related Work}

Open reasoning models have been advanced primarily through post-training. Reinforcement learning with
verifiable rewards traces to outcome-based reward modeling \cite{cobbe2021gsm8k} and was prominently
applied to mathematical reasoning \cite{lightman2023lets, shao2024deepseekmath}; our GRPO usage
follows the group-relative policy-gradient approach. Long chain-of-thought supervision builds on
instruction tuning \cite{wei2022finetuned} and preference optimization \cite{ouyang2022instructgpt,
rafailov2023dpo}. The base model, Qwen3-8B \cite{qwen3report}, already incorporates strong reasoning
behavior via the Qwen team's own post-training; Zen-Pro adapts and repackages this base rather than
competing with larger models on parameter count.

%% ─────────────────────────────────────────────────────────────────────────────
\section{Licensing and Attribution}

Qwen3-8B is released by Alibaba Cloud under the Apache-2.0 license, permitting commercial use,
modification, and redistribution subject to attribution. Zen-Pro is distributed under the same
Apache-2.0 terms with upstream attribution and NOTICE. Authoritative terms are those of the upstream
Qwen3-8B model card and LICENSE.

%% ─────────────────────────────────────────────────────────────────────────────
\section{Limitations}

Zen-Pro inherits the limitations of Qwen3-8B. As an 8B model, its reasoning depth and world
knowledge are bounded by that scale; reasoning-focused fine-tuning improves chain-of-thought
formatting and, on favorable domains, accuracy, but does not match the capability of substantially
larger models. RLVR helps only where correctness is programmatically verifiable; commonsense and
open-ended tasks rely on the base model and SFT signal. The model can still hallucinate, and
alignment is incomplete. Fine-tuning can also regress some base-model capabilities; any claimed
improvement should be verified against the Qwen3-8B baseline.

%% ─────────────────────────────────────────────────────────────────────────────
\section{Conclusion}

Zen-Pro is an honest, reasoning-tuned, well-packaged derivative of the Apache-2.0 Qwen3-8B model---an
8B model, not a 72B one, and not trained from scratch. Its contribution is reasoning-focused
post-training (long-CoT SFT and optional RLVR) and convenient deployment artifacts. We document the
inherited 8B architecture transparently and defer to the upstream Qwen3 technical report for
pretraining details and rigorous benchmark numbers.

%% ─────────────────────────────────────────────────────────────────────────────
\begin{thebibliography}{99}

\bibitem{qwen3report}
Qwen Team, Alibaba Cloud, ``Qwen3 Technical Report,'' \textit{arXiv:2505.09388}, 2025.

\bibitem{qwen3hf}
Qwen Team, ``Qwen3-8B model card and configuration,'' Hugging Face,
\url{https://huggingface.co/Qwen/Qwen3-8B}, 2025.

\bibitem{lightman2023lets}
H.~Lightman et al., ``Let's Verify Step by Step,'' \textit{arXiv:2305.20050}, 2023.

\bibitem{cobbe2021gsm8k}
K.~Cobbe et al., ``Training Verifiers to Solve Math Word Problems,'' \textit{arXiv:2110.14168}, 2021.

\bibitem{shao2024deepseekmath}
Z.~Shao et al., ``DeepSeekMath: Pushing the Limits of Mathematical Reasoning in Open Language
Models,'' \textit{arXiv:2402.03300}, 2024.

\bibitem{peng2023yarn}
B.~Peng et al., ``YaRN: Efficient Context Window Extension of Large Language Models,''
\textit{arXiv:2309.00071}, 2023.

\bibitem{ainslie2023gqa}
J.~Ainslie et al., ``GQA: Training Generalized Multi-Query Transformer Models,''
\textit{EMNLP}, 2023.

\bibitem{wei2022finetuned}
J.~Wei et al., ``Finetuned Language Models Are Zero-Shot Learners,'' \textit{ICLR}, 2022.

\bibitem{rafailov2023dpo}
R.~Rafailov et al., ``Direct Preference Optimization,'' \textit{NeurIPS}, 2023.

\bibitem{ouyang2022instructgpt}
L.~Ouyang et al., ``Training Language Models to Follow Instructions with Human Feedback,''
\textit{NeurIPS}, 2022.

\end{thebibliography}

\end{document}
